\subsection{Network and Memory Architecture}
\label{app:arch_details}

Each memory injection block of Sec.~\ref{sec:memmethod:integration}
(Fig.~\ref{fig:bridgevla_arch}) stacks $L{=}2$ layers with 8 attention heads of dimension 128 and a
feed-forward expansion factor of 2, and operates in the
2048-dimensional patch-token space of the backbone.
The temporal memory uses two such blocks, one for the initial anchor
views and one for the dynamic keyframe bank, and the spatial memory a
third (Fig.~\ref{fig:anchor_memory}); the anchor block concatenates the three views so that attention
can track scene changes across views, whereas the other two blocks
restrict each view's tokens to the corresponding memory view.

\subsection{Pre-Training}
\label{app:pretrain_details}

All fine-tuning runs in this article warm-start from a single
2D-heatmap pre-training run that instantiates
Sec.~\ref{sec:bridgevla:pretrain} on the 120K object-detection split
of RoboPoint~\cite{yuan2024robopoint}; Fig.~\ref{fig:pretrain_dataset}
illustrates how the ground-truth heatmaps are rendered from the
detection boxes, and Fig.~\ref{fig:real_heatmaps} the predictions the
fine-tuned model still produces on such data.
The memory injection blocks and the sub-goal gate of
Sec.~\ref{sec:memmethod} are not pre-trained and are instead trained
from scratch during fine-tuning.
Whenever a stage has nothing to read, the injection block gates all
residual contributions of masked memory to zero, so the memory-free
\method\ forward pass is recovered exactly.

\subsection{Fine-Tuning Details}
\label{app:finetune_hparams}

\begin{table*}[t]
\centering
\caption{\textbf{Per-benchmark fine-tuning configuration.}
Fine-tuning is two-phase: during the initial freeze epochs the
PaliGemma backbone is frozen and only the modules outside it are
trained (the convex-upsampling modules, the action heads, and the
memory modules); the backbone is then unfrozen for the remaining
epochs, and the learning-rate warmup applies per phase.}
\label{tab:hyperparams}
\footnotesize
\setlength{\tabcolsep}{6pt}
\setlength{\aboverulesep}{0pt}
\setlength{\belowrulesep}{0pt}
\renewcommand{\arraystretch}{1.2}
\begin{tabular}{lccccc}
\toprule
\textbf{Setting} & \textbf{RLBench} & \textbf{COLOSSEUM} & \textbf{GemBench} & \textbf{RMBench} & \textbf{MemoryBench} \\
\midrule
Epochs / freeze epochs & 130 / 4 & 200 / 4 & 200 / 2 & $\sim$320 (per task) / 5 & 160 / 20 \\
Warmup steps (per phase) & 1,500 & 1,500 & 1,000 & 1,000 & 1,000 \\
Optimizer & \multicolumn{5}{c}{AdamW, $(\beta_1, \beta_2) = (0.9, 0.95)$} \\
Learning rate & $8{\times}10^{-5}$ & $8{\times}10^{-5}$ & $5{\times}10^{-5}$ & $5{\times}10^{-5}$ & $5{\times}10^{-5}$ \\
Weight decay & $10^{-2}$ & $10^{-2}$ & $10^{-3}$ & $10^{-3}$ & $10^{-3}$ \\
Batch size per GPU & 4 & 4 & 4 & 4 & 4 \\
GPUs & 32 & 32 & 32 & 8 & 8 \\
\midrule
SE(3) aug.\ (trans.\ / yaw) & 0.125 / $45^{\circ}$ & 0.125 / $45^{\circ}$ & 0.125 / $45^{\circ}$ & 0.03 / $30^{\circ}$ & 0.125 / $45^{\circ}$ \\
Stage-2 zoom jitter & 0.05 & 0.05 & 0.05 & 0.005 & 0.05 \\
\midrule
Slot budget $K$ & 2 & 2 & 2 & 12 & 2 \\
Neighboring keyframes $n$ & 2 & 2 & 2 & 2 & 2 \\
Gate $\lambda_{\mathrm{check}}$ / pos.\ wt.\ / thr. & -- & -- & -- & 1.0 / 5.5 / 0.5 & -- \\
\midrule
Rotation head & \multicolumn{5}{c}{continuous 6D~\cite{zhou2019continuity}} \\
Collision head & on & on & off & off & off \\
Arms & 1 & 1 & 1 & 2 & 1 \\
Demonstrations per task & 100 & 100 & 100 (per variation) & 50 & 100 \\
\midrule
Input cameras (resolution) & 4 ($128^2$) & 4 ($128^2$) & 4 ($256^2$) & 4 ($224^2$) & 4 ($128^2$) \\
Ortho.\ window scale & 2.0 & 2.0 & 2.0 & 0.8 & 2.0 \\
Splat radius (coarse / fine) & 0.012 / 0.012 & 0.012 / 0.012 & 0.012 / 0.012 & 0.004 / 0.012 & 0.012 / 0.012 \\
\bottomrule
\end{tabular}
\end{table*}

Table~\ref{tab:hyperparams} lists the per-benchmark fine-tuning
configuration; the paragraphs below cover the settings the table
cannot express.

\paragraph{Two-phase schedule}
During the initial freeze epochs of Table~\ref{tab:hyperparams}, the
PaliGemma backbone is frozen and gradients reach only the modules
outside it: the convex-upsampling modules, the MLP action heads, and
the memory injection blocks and sub-goal gate.
This first phase lets the modules that the pre-training of
Appendix~\ref{app:pretrain_details} does not cover adapt to the
manipulation data before the backbone is touched.
In the second phase the backbone is unfrozen and trained jointly with
these modules, except for the SigLIP vision encoder and the
language-token embedding, which remain frozen throughout fine-tuning.
The optimizer is re-initialized at the phase boundary, so the warmup
steps of Table~\ref{tab:hyperparams} apply to each phase.

\paragraph{Memory-specific settings}
Whenever memory is enabled, random in-plane 2D image augmentation is
disabled and the workspace cube is centered on fixed scene bounds
rather than on the per-frame cloud mean, since either perturbation
would break the pixel correspondence between the current observation
and the cached memory tokens.
The SE(3) augmentation of Table~\ref{tab:hyperparams}, whose
translation is a fraction of the workspace extent and whose rotation
perturbs yaw only, is instead applied jointly to the current, anchor,
and history point clouds (Sec.~\ref{sec:memmethod:training}); the
stage-2 zoom jitter perturbs the ground-truth waypoint the fine stage
zooms to during training.

\paragraph{Slot budgets}
On RMBench the budget $K{=}12$ holds the two neighboring keyframes and
up to ten sub-goal slots.
Every executed keyframe occupies a neighboring slot regardless of the
gate, and a gated keyframe enters a sub-goal slot only when it leaves
the neighboring window two steps later; when the sub-goal slots are
full, the oldest is evicted.
The other benchmarks carry no sub-goal annotations, so the gate is
disabled and the budget holds only the two neighboring keyframes,
$K{=}2$; the temporal memory there reduces to the neighboring
keyframes and the initial anchor, which suits the shorter horizons of
these tasks.

\paragraph{Rendering}
Every point cloud is rendered into three $224 \times 224$ orthographic
views regardless of the sensor resolution listed in
Table~\ref{tab:hyperparams}.
The orthographic window scale sets the extent of the rendered viewport
relative to the workspace, with values below one zooming in.

\paragraph{Ablation configurations}
Among the design-ablation rows of Table~\ref{tab:rlbench}
(Sec.~\ref{sec:exp:ablation}), \emph{w/ discretized rotation} replaces
the continuous 6D rotation head with per-axis Euler angles quantized
into $5^{\circ}$ bins and supervised with cross-entropy, as in our
conference version.
Beyond its resolution ceiling, this representation is ill-conditioned
near the gimbal-lock singularities of the Euler decomposition, where
the roll and yaw axes degenerate and orientations that are close in
$SO(3)$ may fall into distant bins, making the per-axis cross-entropy
targets discontinuous.
The continuous 6D parameterization is free of such singularities.
\emph{w/ 3D position input} adds a 3D convolutional module that encodes
per-pixel 3D positions and fuses them with the 2D image features fed to
the backbone.
\emph{w/o heatmap decoding} replaces the convex-upsampling module, of
309M parameters, with a similarly sized Transformer decoder of 303M
parameters that regresses the target positions directly under an MSE
loss, leaving all other modules unchanged.

\subsection{Training Data Preparation}
\label{app:keyframe_selection}

\paragraph{Keyframe selection}
Demonstrations are converted into consecutive-keyframe training
transitions with the keyframe-selection strategy of
PerAct~\cite{shridhar2023perceiver} in all single-arm experiments: a
time step is labeled a keyframe if the robot is stationary, if the
gripper state changes, or if it is the final step of the episode.
RMBench demonstrations instead use a bimanual variant of this
heuristic, which keeps the last frame of every segment in which both
arms are still.

\paragraph{Sub-goal labels}
\label{app:gate_labels}
The sub-goal gate of Sec.~\ref{sec:memmethod:gate} is supervised by
labels derived from the RMBench demonstrations, whose keyframes carry
per-segment language annotations.
The last keyframe of each language segment marks the completion of a
sub-goal and is labeled positive; all other keyframes are negatives.
Repeated identical language segments are kept separate rather than
merged by text identity, so every repetition of an instruction, such
as each press in \emph{Press Button}, contributes its own sub-goal
frame.
The resulting positives cover 9--15\% of keyframes depending on the
task, which the binary cross-entropy compensates with the
positive-class weight of Table~\ref{tab:hyperparams}.

\subsection{Evaluation Protocol}
\label{app:eval_protocol}

This appendix states the training data, trial counts, step budgets,
and reporting statistics used for each benchmark.

\paragraph{RLBench}
Training uses the 100 demonstrations per task provided by the
benchmark, over the 18 tasks of Fig.~\ref{fig:RLBench_VIS}.
Each configuration is evaluated over 25 episodes per task under a
25-keyframe-step budget, except \emph{Place Cups} and \emph{Stack
Blocks}, which receive 35 steps.
We report means and standard deviations over five independent
evaluation runs, except for the two design-ablation variants that
replace the heatmap head or inject 3D position input, which are
evaluated over three runs.

\paragraph{COLOSSEUM}
Policies are trained on the unperturbed RLBench data of the 20
benchmark tasks, 100 demonstrations per task, and evaluated under the
14 settings of Table~\ref{tab:colosseum}, visualized in
Fig.~\ref{fig:vis_colosseum}, with 25 trials per task and setting.
\method, \memmethod, and RVT-2 are our own training and evaluation
runs, reported as mean and variance over three test repetitions
(Tables~\ref{tab:results_bridgevla},
\ref{tab:results_bridgevla_plus},
and~\ref{tab:results_rvt2}); the remaining baseline numbers are quoted
from the benchmark release~\cite{pumacay2024colosseum}.

\paragraph{GemBench}
Policies are trained on the 16-task training split, 31 variations, and
evaluated on the 44 test tasks, 92 variations, of the four levels
L1--L4 shown in Fig.~\ref{fig:vis_gembench}.
Following the benchmark protocol~\cite{garcia2024towards}, \method\ and
\memmethod\ are each evaluated over five random seeds with 20 trials
per task variation, and Tables~\ref{tab:gembench}
and~\ref{tab:gembench_sota_cmpr_l1_detail}--\ref{tab:gembench_sota_cmpr_l4_detail}
report means over the five seeds.
Baseline numbers are quoted from~\cite{garcia2024towards}.

\paragraph{RMBench}
The benchmark is built within RoboTwin~2.0~\cite{chen2025robotwin2} and
simulated in SAPIEN~\cite{xiang2020sapien}; its nine dual-arm tasks are
shown in Fig.~\ref{fig:vis_rmbench}.
Following the benchmark protocol, policies are trained on 50 expert
demonstrations per task; \memmethod\ emits one action tuple per arm at
every keyframe step (Sec.~\ref{sec:bridgevla:problem}), and the
collision flag is dropped.
Every reported number is a single 100-episode evaluation under a
per-task keyframe-step limit scaled to the task horizon.
The numbers in Tables~\ref{tab:rmbench_main}
and~\ref{tab:rmbench_ablation} are obtained by training one model per
task and selecting the best-performing checkpoint of each task's training run.

\paragraph{MemoryBench}
Policies are trained on 100 demonstrations per task following the
protocol of~\cite{fang2025sam2act}.
One evaluation covers all nine task variants of the three tasks of
Fig.~\ref{fig:vis_memorybench} under a 25-step budget, and we report mean$\pm$std over five evaluation seeds.

\paragraph{Real robot}
Each of the 13 Franka tasks of Figs.~\ref{fig:basic_task1}
and~\ref{fig:basic_task2} is trained with 10 kinesthetic-teaching
demonstrations, reduced to 3 in the low-data variant of
Table~\ref{tab:app:real_episode}, and every method is evaluated over 10
trials per task in the Basic setting.
Each test scene is photographed and manually aligned across methods.
A single multi-task model is trained jointly on all tasks, and the
same checkpoint is evaluated on every task.
The six generalization settings compare \method\ and RVT-2, the two
methods that perform well in the Basic setting; their definitions are
given in Appendix~\ref{app:real_settings_vis} and the four
visual-disturbance ones are visualized in Fig.~\ref{fig:real_settings}.

The Dobot suite of Appendix~\ref{app:real_dobot} follows the same
protocol at the granularity of the language instruction rather than the
task: its five tasks, three memory-dependent and two memory-free,
comprise seven instructions, each trained with 10 kinesthetic-teaching
demonstrations, for 70 demonstrations in total.
All methods share this training set, and \method\ and \memmethod\ are
each trained jointly on all seven instructions as a single model whose
one checkpoint is evaluated on every instruction, as on the Franka
platform.
Every instruction is evaluated
over 10 trials in each of the Basic, Distractor, Background, Height,
and Lighting settings, the last four illustrated in
Fig.~\ref{fig:dobot_settings} and defined as in the Franka suite
(Appendix~\ref{app:real_settings_vis}).
The comparison here runs across all five
settings and against two baselines, the memory-free \method\ and the
memory-augmented SAM2Act+~\cite{fang2025sam2act}, so that the effect of
memory is separated from that of the backbone: \method\ shares
\memmethod's backbone but has no memory, whereas SAM2Act+ has memory
but a different backbone and retrieval scheme.
Figs.~\ref{fig:dobot_rollouts1} and~\ref{fig:dobot_rollouts2} show
\memmethod\ rollouts on the memory-dependent and memory-free
instructions, respectively.

\subsection{Computational Cost}
\label{app:compute_resources}

The running times below are those of the runs behind the reported
results.
The 2D-heatmap pre-training of Appendix~\ref{app:pretrain_details}
takes about 2 hours on 8 NVIDIA A100 GPUs.
For fine-tuning, \method\ and \memmethod\ train on 32 H20 GPUs on
RLBench, COLOSSEUM, and GemBench, as do the ablation rows of
Table~\ref{tab:rlbench}; RMBench trains one model per task on 8 H20
GPUs each, and MemoryBench trains on 8 A100 GPUs.
Real-world fine-tuning takes about 1.5 hours on 8 A100 GPUs.

Evaluation uses a single GPU per run: RMBench, GemBench, MemoryBench,
and COLOSSEUM are evaluated on one A100, and RLBench on one H20.
For real-world deployment, both models run on a machine with a single
NVIDIA RTX 4090 GPU.
Averaged over 100 trials, the end-to-end latency from point-cloud input
to action output is 0.35 seconds per prediction step for \method\ and
0.57 seconds for the full \memmethod\
(Sec.~\ref{sec:memmethod:training}).
Both figures are small relative to the remainder of the control loop:
in real-world deployment, transmitting the multi-camera observations
and physically executing the planned motion between keyframes dominate
the total time per step.

\paragraph{Memory overhead}
The totals of Sec.~\ref{sec:memmethod:training} decompose as follows.
Each of the three memory injection blocks holds 83.95M parameters and
the sub-goal gate a further 17.91M, for 269.77M in total, or 9.2\% of
the 2.92B-parameter backbone.
A cached memory entry is a single bf16 coarse-stage token grid of
3.0\,MiB, so a full memory at the RMBench budget of $K{=}12$ frames
plus the anchor occupies 39\,MiB.
Because entries are stored already encoded, the cache is the only
stored quantity that grows with $K$; the cross-attention cost of the
injection block also grows linearly in $K$ but remains negligible next
to a backbone forward pass, and the number of forward passes per step
is independent of $K$.
Per step, the memory-free policy runs two backbone forwards, coarse and
fine, and the full \memmethod\ three, the additional pass encoding the
fine-stage geometric reference.
Dual-arm deployments, whose fine stage runs once per arm on a shared
coarse trunk (Sec.~\ref{sec:memmethod:bimanual}), accordingly run five.
